\documentclass[tikz,border=10pt]{standalone}
\usepackage{pgfplots}
\usetikzlibrary{calc}

\begin{document}
\begin{tikzpicture}[
    line join=round, 
    line cap=round, 
    >=stealth, 
    thick,
    dot/.style={circle, fill=black, inner sep=1pt}
]

    % --- 1. 定义坐标点 ---
    % 设定 A 为原点，B 为水平轴上一点，C 位于上方，P 为 AC 的延长线上
    \coordinate (A) at (0,0);
    \coordinate (B) at (4,0);
    \coordinate (C) at (5,2);
    
    % 使用 calc 库计算 P 点，使 P 在 AC 的延长线上
    % 这里设定 P 的长度为 AC 长度的 1.5 倍
    \coordinate (P) at ($(A)!1.5!(C)$);

    % --- 2. 绘制图形 ---
    % 绘制三角形 ABC
    \draw (A) -- (B) -- (C) -- cycle;
    
    % 绘制射线 CP (实际上是绘出直线段 AP)
    \draw (C) -- (P);

    % --- 3. 标注顶点标签 ---
    % 根据图像位置调整标签方向
    \node[left] at (A) {$A$};
    \node[below] at (B) {$B$};
    \node[right] at (C) {$C$};
    \node[right] at (P) {$P$};

    % --- (可选) 标注点 ---
    % 如果需要显示点的小圆黑点，可以取消下面代码注释
    %\foreach \p in {A,B,C,P} \node[dot] at (\p) {};

\end{tikzpicture}
\end{document}